\textbf{COLUNI 2018 -- Questão 36}\par\medskip Segundo o historiador Eric Hobsbawm, o século XX começa com a eclosão da Primeira Guerra Mundial em 1914. O conflito durou cerca de quatro anos, finalizando com a rendição dos alemães em 1918. Um dos efeitos mais marcantes desse evento foi o revanchismo alemão que em grande parte causou outra guerra: a Segunda Guerra Mundial (1939-1945). Entre as duas grandes guerras, alguns eventos marcaram a história.

\par

Sobre o chamado período entre-guerras (1919-1939), é INCORRETO afirmar que:\begin{enumerate}
\item A Revolução Russa se consolida com a vitória do Exército Vermelho sobre o Exército Branco ao final da guerra
        civil em 1921.
\item O modelo liberal entra em crise, levando à ascensão de lideranças políticas totalitárias e antidemocráticas
        em alguns países da Europa.
\item A crise de 1929, iniciada com o \textit{crack} da bolsa de Nova Iorque nos EUA, gera efeitos desastrosos nas
        demais economias capitalistas do mundo, aumentando a intervenção estatal.
\item A expansão alemã nazista começa com a anexação pacífica da Polônia, em 1938, e a imediata ruptura do
        pacto de neutralidade firmado com a URSS.\end{enumerate}
